% Further C.1 deductions; owned by cycles. No preamble.
\paragraph{C.1j (Proved: the exact middle-positivity obstruction over $\mathbb C$).}
Let $X/\mathbb C$ be smooth projective of dimension $2m$, and put
\[
 P^m_{\mathrm{alg},L}(X)
 =\ker\bigl(L:A^m_{\mathrm{alg}}(X)\longrightarrow
                         A^{m+1}_{\mathrm{alg}}(X)\bigr).
\]
The form $q(x,y)=(-1)^m\deg(xy)$ is positive semidefinite on this
space. Its radical is exactly
\begin{equation}\label{eq:Balg-exact-gate}
 \operatorname{rad}(q)=
 \ker\bigl(L:G^m(X)\longrightarrow G^{m+1}(X)\bigr).
\end{equation}
Consequently the letter's $B(X,L)$ is equivalent to the injectivity
of this particular map on the rational Griffiths group.

\emph{Proof.}
The cycle class of $x\in P^m_{\mathrm{alg},L}(X)$ is a primitive
rational Hodge class of type $(m,m)$. The Hodge--Riemann bilinear
relations say that $(-1)^m\int_X\mathrm{cl}(x)^2$ is positive
unless $\mathrm{cl}(x)=0$; this is precisely the complex positivity
input recorded in \citet[Remark 3.10, p.~380]{cyKleiman1968}.
The intersection number agrees with that integral. Thus $q(x,x)\ge0$,
and it is zero exactly for $x\in G^m(X)$. Such an $x$ has intersection
zero with every algebraic cycle, so it is in the radical. Conversely
a radical vector has $q(x,x)=0$, hence has zero cycle class. Taking
the intersection with the defining primitive kernel proves
\eqref{eq:Balg-exact-gate}.

This proves $B$ for $m=0,1$, for every $X$ with $G^m(X)=0$, and more
generally whenever the displayed Griffiths map is injective. In all
complex dimensions the weaker nonnegative form is already a theorem.
Our counterexamples to Lefschetz \emph{surjectivity} do not supply a
nonzero vector in \eqref{eq:Balg-exact-gate}. We therefore leave the
universal strict-positivity assertion for $m\ge2$ undecided, rather
than infer its failure from the existence of a Griffiths group.

\paragraph{C.1k (Proved: exact universal ranges for three consequences).}
Fix $n\ge0$ and $0\le r\le n$. Over $\mathbb C$, each of the following
properties holds for every smooth projective connected $n$-fold
exactly when $r\in\{0,1,n\}$, with repeated or inadmissible entries
omitted:
\begin{enumerate}
\item the map $A^r_{\mathrm{alg}}(X)\to H^{2r}(X,\mathbb Q)$ is
injective;
\item algebraic equivalence up to torsion equals numerical equivalence
in codimension $r$;
\item $A^r_{\mathrm{alg}}(X)$ is finite dimensional over $\mathbb Q$.
\end{enumerate}
The corresponding assertion that cycles modulo algebraic equivalence
up to torsion form a finitely generated abelian group has the same
universal range.

For $r=0,n$, degree gives the assertions. For $r=1$, the group is
the N\'eron--Severi group modulo its finite torsion, which is finitely
generated. The divisor cycle map is injective after tensoring with
$\mathbb Q$; the Hodge index theorem identifies numerical and
homological equivalence for divisors. These are the usual
N\'eron--Severi and Lefschetz $(1,1)$ statements, also used in
C.1c--d and \citet[\S3, especially 3.4 and 3.10]{cyKleiman1968}.

Every remaining pair has $n\ge3$ and $2\le r\le n-1$. Take the
very general quintic $V$ of C.1a and
$X=V\times\mathbb P^{n-3}$. The projective-bundle formula supplies
the direct summand
\[
 G^2(V)t^{r-2}\ \subset\ G^r(X),
 \qquad 0\le r-2\le n-3.
\]
It is infinite dimensional by the opened Clemens theorem. Every
element is numerically trivial, and a nonzero one disproves (1) and
(2); the infinite-dimensional summand disproves (3) and finite
generation. This proof concerns universal assertions. For a specified
variety the first property is exactly $G^r(X)=0$, while equality with
numerical equivalence may additionally require the homological
standard conjecture $D$.

The letter's image criterion also fails in each of these target
codimensions, even for the zero correspondence: its cohomological
image contains the class of every $g\in G^r(X)$, but its cycle image
does not contain a nonzero $g$. More generally, that criterion for
all correspondences into a fixed $A^r_{\mathrm{alg}}(X)$ necessarily
implies $G^r(X)=0$. This necessary condition alone does not prove
the criterion, since an arbitrary cohomological preimage must still
be represented by an algebraic cycle.

\paragraph{C.1l (Proved: a nonzero radical in the correspondence algebra).}
If $G^r(X)\ne0$ for some $r$, the algebra
\[
 R=A^{n}_{\mathrm{alg}}(X\times X),\qquad n=\dim X,
\]
with composition of degree-zero correspondences, has nonzero
Jacobson radical. In particular it is not a semisimple algebra.
This supplies counterexamples to the letter's correspondence-ring
consequence in every dimension $n\ge3$.

Choose $0\ne g\in G^r(X)$, an ample class $h$, and
$d=\deg h^n>0$. Define
\[
 a=d^{-1}(h^{n-r}\times g)\in R.
\]
It acts by
$a_*(x)=d^{-1}\deg(xh^{n-r})g$ on $A^r_{\mathrm{alg}}(X)$.
In particular $a_*(h^r)=g$, so $a\ne0$ as a correspondence class.
For every $b\in R$, the usual composition formula for decomposable
correspondences gives
\[
 aba=d^{-1}\deg\bigl(h^{n-r}\,b_*g\bigr)a=0.
\]
Indeed $b_*g$ is homologically trivial, hence numerically trivial.
It follows that $(ba)^2=0$, so $1-ba$ is invertible, with inverse
$1+ba$, for every $b\in R$. The defining radical criterion for a
unital ring now gives $a\in\operatorname{Jac}(R)$. This proves
nonsemisimplicity; merely exhibiting a nilpotent element would not
have sufficed, since matrix algebras themselves contain nilpotents.
The product examples in C.1k finish the claimed dimension range.

There is also a direct obstruction to a semisimple category of
motives modulo algebraic equivalence. A nonzero Griffiths cycle
gives a nonzero morphism from a Tate motive to the appropriate
twist of $h(X)$, whereas composition with every morphism back is
zero by numerical triviality. In a semisimple category with
endomorphism ring $\mathbb Q$ for that Tate object, every nonzero
such morphism would split. Thus this category cannot be semisimple.
These arguments neither refute semisimplicity of numerical motives
nor claim semisimplicity for an individual variety without checking
its correspondences on the product. C.1m supplies the missing
universal surface case.

\paragraph{C.1m (Proved under an opened classical input: the exact universal
dimension range for finite semisimple correspondence algebras).}
For smooth projective connected complex $n$-folds, the assertion
that $A^n_{\mathrm{alg}}(X\times X)$ is a finite semisimple
$\mathbb Q$-algebra holds universally exactly for $n\le1$.

For $n=0$ the algebra is $\mathbb Q$. For a smooth projective
curve $C$, normalized divisor correspondences and a chosen point
give an algebra isomorphism
\[
 A^1_{\mathrm{alg}}(C\times C)
    \simeq\mathbb Q\times\operatorname{End}(J(C))_{\mathbb Q}
                        \times\mathbb Q.
\]
The three factors act on $H^0,H^1,H^2$, respectively. The
divisor cycle map is injective, so there are no further
off-diagonal components. The middle identification is the
classical Picard-correspondence identification used in C.4a.
Poincar\'e complete reducibility makes
$\operatorname{End}(J(C))_{\mathbb Q}$ a finite product of matrix
algebras over division algebras. Hence the displayed algebra
is finite and semisimple.

For a specified surface counterexample, let the fourteen
$t_\alpha$, indexed by $|\alpha|=4$ and
$\alpha\ne(4,0,0)$, be algebraically independent over
$\mathbb Q$, and put
\[
 K=\mathbb Q(t_\alpha),\qquad
 F=x_0^4+\sum_{\substack{|\alpha|=4\\\alpha\ne(4,0,0)}}
                      t_\alpha x^\alpha.
\]
Let $C\subset\mathbb P^2_{\overline K}$ be $F=0$.
It is smooth, since the discriminant is not identically zero
on this coefficient chart, as the Fermat quartic shows.
Choose a root $\rho\in\overline K$ of $F(T,1,0)$ and put
$p=(\rho:1:0)$. Embed $\overline K$ into $\mathbb C$.
We claim that the modified diagonal
\[
 \gamma=\Delta_{123}-\Delta_{12}-\Delta_{13}-\Delta_{23}
                     +\Delta_1+\Delta_2+\Delta_3
             \quad\text{in }\mathrm{CH}^{2}(C^3)
\]
has infinite order modulo algebraic equivalence. Here the
coordinates indexed by a subscript vary together and the other
coordinates equal $p$. The opened primary source
\citet[Proposition 3.1, p.~654, and the paragraph immediately
after Corollary 4.9, p.~659]{cyGrossSchoen1995} proves its
homological triviality and records the existence of such
genus-three curves, following the Ceresa and Colombo--van Geemen
results. These examples are nonhyperelliptic: the hyperelliptic
modified diagonal is torsion by their Corollary 4.9.
Thus they occur among smooth plane quartics. We use that
classical nonvanishing input explicitly;
we do not claim an independent proof of Ceresa's theorem or
that its original Annals PDF was obtained. The formula for
$\gamma$ is equation (0.1) of the same opened paper.

We justify the passage from that existence input to this
specified geometric generic curve. The class of $\gamma$
modulo algebraic equivalence is independent of $p$, since
the displayed seven-term formula varies algebraically as
$p$ ranges over the connected curve. Suppose that a nonzero
multiple of the geometric generic $\gamma$ were algebraically
trivial. Its finite cycle, parameter-curve, and rational-function
data descend to a finite extension of the generic coefficient
field. Norm and spreading then make a nonzero integer multiple
algebraically trivial over a nonempty open of the smooth
plane-quartic parameter space. Such a relation specializes to
every smooth fibre: choose an integral curve through the
desired parameter point and meeting that open, normalize,
make the finite base change needed for the witnesses, and
specialize at a point over the desired fibre. Algebraic
equivalence is preserved in this specialization. Concretely,
one closes the witnessing parameter curves and correspondences
over the resulting discrete valuation ring; the proper
connected special parameter curves, component by component,
give the specialized algebraic-equivalence relation. Multiplying
by the finite base-change degree has no effect on the
contradiction with infinite order. Every smooth plane quartic
can be put in the displayed coefficient chart by choosing
coordinates with nonzero $x_0^4$ coefficient. We would therefore
force the cited nonvanishing example to be torsion, a contradiction.
This proves the asserted non-torsion class for the specified
$C$, for any chosen embedding into $\mathbb C$; extension of
algebraically closed fields does not create a new finite
algebraic-equivalence relation, by the same finite-data
specialization argument.

Put $S=C\times C$, and view $\gamma$ as a degree-zero
correspondence $S\to C$. If $\pi:S\to C$ is the first
projection, put $a=\pi^*\circ\gamma$. In coordinates on
$S\times S=C^4$ this is the class $\gamma\times C$.
Pullback to a fixed point in the fourth factor recovers
$\gamma$, so $a\ne0$ modulo algebraic equivalence with
rational coefficients.
For every $b\in A^2_{\mathrm{alg}}(S\times S)$, the composition
\[
       \gamma\circ b\circ\pi^*
          \in A^1_{\mathrm{alg}}(C\times C)
\]
is homologically trivial. It is a divisor class, hence zero
modulo algebraic equivalence. Consequently
\[
 aba=\pi^*(\gamma b\pi^*)\gamma=0
                         \quad\text{for every }b.
\]
The radical argument in C.1l proves $0\ne a\in\operatorname{Jac}(R_S)$.
No nilpotence conjecture or finite-dimensional-motive hypothesis
is used. In particular this surface satisfies both of the
fixed-variety assertions $A(S)$ and $B(S)$, while its
correspondence algebra is not semisimple. The universal
auxiliary-product hypotheses in the letter's motivic argument
cannot be replaced by hypotheses on $S$ alone.

Finally take $X=S\times\mathbb P^{n-2}$ for $n\ge2$, with
projection $f:X\to S$ and a constant section $s:S\to X$.
The class $a_X=f^*a s^*$ is nonzero because
$s^*a_X f^*=a$. For every $b\in R_X$,
\[
 a_X b a_X=f^*a(s^*b f^*)a s^*=0.
\]
Thus it also lies in the nonzero Jacobson radical. This gives
the counterexample in every dimension $n\ge2$ and proves
the stated universal range. Notice that the Griffiths groups
of the surface $S$ itself vanish; it is the Griffiths group
of $S\times S$ that is detected here.
